\documentclass[11pt,twoside]{article}
\usepackage[T1]{fontenc}
\usepackage[latin1]{inputenc}
\usepackage[english]{babel}
\usepackage{amsmath}
\usepackage{amscd}
\usepackage{amssymb}
\usepackage{multirow}
\usepackage{tabularx}
\usepackage{url}
\usepackage{fancyhdr}
\usepackage{lastpage}
\usepackage{lipsum}
\usepackage{hyperref}
\usepackage[a4paper,margin=3cm,hmarginratio=1:1]{geometry}
\usepackage{wasysym}

%%%%%%%%%%%%%%%%%%%%%%%%%%%%%%%%%%%%%%%%%%%%%%%%%%%%%%%%%%%%%%%%%%%%%%%%%%%
%%%%%%%%%%%%% DO NOT TOUCH ANYTHING BELOW THIS LINE %%%%%%%%%%%%%%%%%%%%%%%
%%%%%%%%%%%%%%%%%%%%%%%%%%%%%%%%%%%%%%%%%%%%%%%%%%%%%%%%%%%%%%%%%%%%%%%%%%%

\newcommand{\surveyproject}{Survey Project}
\newcommand{\coursenumber}{DD2448}
\newcommand{\coursename}{\coursenumber~Foundations of cryptography}
\newcommand{\coursenick}{krypto26}

\lhead{\coursename~(\coursenick)}
\chead{}
\rhead[\coursename~(\coursenick)]{}
\lfoot[\thepage~(\pageref{LastPage})]{}
\cfoot{}
\rfoot[]{\thepage~(\pageref{LastPage})}

\fancypagestyle{firststyle}
{
   \fancyhf{}
   \fancyfoot[R]{\thepage~(\pageref{LastPage})}
}

\renewcommand{\headrulewidth}{0pt}

%%%%%%%%%%%%%%%%%%%%%%%%%%%%%%%%%%%%%%%%%%%%%%%%%%%%%%%%%%%%%%%%%%%%%%%%%%%
%%% HERE YOU CAN ADD YOUR OWN MACROS AND ENVIRONMENTS IN THE PREAMBLE %%%%%
%%%%%%%%%%%%%%%%%%%%%%%%%%%%%%%%%%%%%%%%%%%%%%%%%%%%%%%%%%%%%%%%%%%%%%%%%%%

% Add your macros here.

\begin{document}

%%%%%%%%%%%%%%%%%%%%%%%%%%%%%%%%%%%%%%%%%%%%%%%%%%%%%%%%%%%%%%%%%%%%%%%%%%%
%%%%%%%%%%%% THE FOLLOWING GENERATES THE HEADER %%%%%%%%%%%%%%%%%%%%%%%%%%%
%%%%%%%%%%%% DO NOT TOUCH THIS %%%%%%%%%%%%%%%%%%%%%%%%%%%%%%%%%%%%%%%%%%%%
%%%%%%%%%%%%%%%%%%%%%%%%%%%%%%%%%%%%%%%%%%%%%%%%%%%%%%%%%%%%%%%%%%%%%%%%%%%

\newcommand{\gptitle}[1]{
  \begin{center}
    {\LARGE\textbf{#1}}\\
    \medskip
    \today
  \end{center}

  \thispagestyle{fancy}
}

%%%%%%%%%%%%%%%%%%%%%%%%%%%%%%%%%%%%%%%%%%%%%%%%%%%%%%%%%%%%%%%%%%%%%%%%%%%
%%%%%%%%%%%%%%%%%%%%% YOUR SOLUTION STARTS HERE %%%%%%%%%%%%%%%%%%%%%%%%%%%
%%%%%%%%%%%%%%%%%%%%%%%%%%%%%%%%%%%%%%%%%%%%%%%%%%%%%%%%%%%%%%%%%%%%%%%%%%%

\gptitle{You Choose Your Own Title\footnote{Titles should be
    descriptive, preferably short, and catchy if possible.}}

\begin{abstract}
  An abstract should be short and invite the reader to actually read
  the content of your analysis. You may mention what your findings
  are, but you should not say how you reached your conclusions
  here. Including too much information is a common mistake that can be
  found even in reputable researchers' papers. Simple is good.
\end{abstract}

\section{Usually It Is a Good Idea to Introduce the Topic}

\lipsum[2-3]

\subsection{You Choose}

\lipsum[2-3]

\subsection{Your Own Sections and Subsections, But}

\lipsum[2-4]

\subsubsection{Numbered Subsubsections are Possible, and Often a Bad Idea}

\lipsum[2-4]

\subsubsection*{Unnumbered Subsubsections are Often a Better Choice}

\lipsum[2-3]

\paragraph{A paragraph is also possible.}

\lipsum[3-4]

\section{But, You Already Know How to Organize}

\lipsum[2-3]

\section{A Written Report}

\lipsum[2-4]

\section{A Conclusion Is a Nice Ending}

\lipsum[2-3]

\appendix

\section{It Feels Short! What Should We Do?}

We assess quality, not quantity. It may be possible to write a notably
shorter report than the page limit, but we consider this to be
hard. In fact, we expect you to struggle to fit your content within
the page limit.

Focus on content first and keep the page limit vaguely in mind. The
result is probably a little longer than allowed, but squeezing often
improves quality, since you are forced to prioritize.

\section{References Are Entered in a Separate bib-file}

With proper references we can cite Oetiker's excellent introduction to
\LaTeX~\cite{Oetiker2007} or the Swedish election law~\cite{vallag}.

You can also google, e.g., \emph{How to prove yourself DBLP} and find
the DBLP page of Adi Shamir~\cite{dblpshamir}, where you can get the
bibtex code for any of his papers to the left at the arrow symbol,
e.g., the famous paper that introduces the Fiat-Shamir
heuristic~\cite{DBLP:conf/crypto/FiatS86}. Simply copy such entries
into the bib-file and use them.

% The commands below generate a bibliography. You need to run pdflatex
% and bibtex like this to make it work.
%
% For any scientific papers you can usually copy whole bibtex entries
% from ACM or DBLP to your bib-file.
%
% pdflatex SP_solution_template
% bibtex SP_solution_template
% pdflatex SP_solution_template
% pdflatex SP_solution_template
%
% You need to run pdflatex multiple times to ensure that all
% references and citations are correct.
%
\bibliographystyle{plain}
\bibliography{templates/project_solution_template}


\end{document}

%%% Local Variables:
%%% mode: latex
%%% TeX-master: t
%%% End:
